\documentclass[11pt, a4paper]{article}

\usepackage[a4paper, margin=2.5cm]{geometry}
\usepackage{fontspec}
\usepackage[spanish, bidi=basic, provide=*]{babel}
\babelprovide[import, onchar=ids fonts]{spanish}
\babelfont{rm}{Noto Sans}

\usepackage{amsmath}
\usepackage{booktabs}
\usepackage{enumitem}
\setlist[itemize]{label=--}
\usepackage{xcolor}
\usepackage{hyperref}

\definecolor{primary}{RGB}{20, 45, 85}

\begin{document}

% ------------------- PORTADA ACADÉMICA -------------------
\begin{titlepage}
    \centering
    \vspace*{2cm}
    
    {\Huge \bfseries \color{primary} INVESTIGACIÓN DE OPERACIONES}\\[1.5cm]
    
    {\Large \bfseries RESUMEN TÉCNICO Y METODOLÓGICO}\\[2.5cm]
    
    \begin{minipage}{0.8\textwidth}
        \centering
        \large
        \textbf{Maestro:}\\[0.2cm]
        Julio César Galindo López\\[1.5cm]
        
        \textbf{Estudiante:}\\[0.2cm]
        Aranzolo Rios Angel Ariadna\\
    \end{minipage}
    
    \vfill
    
    {\large \today}
    \vspace*{1.5cm}
\end{titlepage}

\newpage
\pagenumbering{arabic}

% ------------------- CONTENIDO DEL RESUMEN -------------------

\section{Formulación de Modelos Lineales}

Un modelo lineal responde a tres preguntas básicas: ¿qué se decide?, ¿qué se quiere lograr? y ¿qué límites no se pueden cruzar?

\subsection*{Receta de 6 pasos}
\begin{enumerate}
    \item Describir la decisión y sus unidades de medida.
    \item Definir una variable de decisión por cada incógnita ($x_1, x_2, \dots, x_n$).
    \item Formular la función objetivo: optimizar $Z = c_1 x_1 + \dots + c_n x_n$ (maximizar beneficio o minimizar costo).
    \item Expresar cada límite como una restricción lineal.
    \item Añadir las condiciones de no negatividad ($x_i \ge 0$).
    \item Comprobar dimensiones, coherencia de unidades y sentido práctico del planteamiento.
\end{enumerate}

\section{Matrices para Organizar Información}

Permiten manejar sistemas grandes de forma compacta mediante la estructura matricial $AX = B$, donde $A$ representa la matriz de coeficientes, $X$ el vector de incógnitas y $B$ el vector de términos independientes.

\subsection*{Operaciones clave}
\begin{itemize}
    \item \textbf{Suma y resta:} Requieren matrices de idénticas dimensiones ($m \times n$).
    \item \textbf{Producto ($AB$):} Exige que el número de columnas de $A$ coincida estrictamente con el número de filas de $B$.
    \item \textbf{Transpuesta ($A^T$):} Intercambia filas por columnas.
    \item \textbf{Resolución de sistemas:} Se estructura la matriz aumentada $[A \mid B]$ y se aplican operaciones elementales de Gauss (intercambiar filas, multiplicar una fila por un escalar no nulo y sumar a una fila el múltiplo de otra).
\end{itemize}

\section{Preparación de un Modelo para Simplex}

El método Simplex trabaja exclusivamente con igualdades, por lo que las desigualdades deben transformarse mediante variables auxiliares con interpretación económica:

\begin{itemize}
    \item \textbf{Variables de holgura ($s \ge 0$):} Se suman en restricciones del tipo $\le$ ($a^T x + s = b$). Representan recursos no utilizados; si su valor final es cero, el recurso asociado se agotó completamente.
    \item \textbf{Variables de exceso ($e \ge 0$):} Se restan en restricciones del tipo $\ge$ ($a^T x - e = b$). Representan cuánto se supera un valor mínimo requerido.
    \item \textbf{Variables artificiales ($a \ge 0$):} Se incorporan en restricciones de igualdad ($=$) o con exceso ($\ge$) para proveer una columna identidad que sirva como base inicial. No representan decisiones reales y deben salir de la base (valer cero) antes de declarar factible la solución.
    \item \textbf{Método de la Gran M:} Penaliza las variables artificiales en la función objetivo asignándoles un coeficiente $M > 0$ sumamente grande ($+M$ en minimización) para obligarlas a abandonar la base durante el proceso iterativo.
\end{itemize}

\section{Rutina y Lectura de la Tabla Simplex}

La tabla se evalúa mediante la fila de costos reducidos ($C_j - Z_j$, donde $Z_j = \sum C_{B_i} a_{ij}$):

\subsection*{Criterios de decisión}
\begin{itemize}
    \item \textbf{Criterio de entrada:}
    \begin{itemize}
        \item En \textbf{maximización}, entra la variable con el mayor costo reducido positivo ($C_j - Z_j > 0$).
        \item En \textbf{minimización}, entra la variable con el costo reducido más negativo ($C_j - Z_j < 0$).
    \end{itemize}
    \item \textbf{Criterio de salida (fila pivote):} Se determina por la razón mínima positiva $\theta = \frac{b_i}{a_{ik}}$, considerando únicamente coeficientes estrictamente positivos en la columna pivote ($a_{ik} > 0$).
\end{itemize}

\subsection*{Actualización (pivoteo)}
\begin{enumerate}
    \item Se normaliza la fila pivote dividiéndola entre el elemento pivote.
    \item Se hacen ceros en las demás filas de la columna pivote aplicando sumas de múltiplos adecuados de la fila pivote.
    \item Se recalculan $Z_j$ y los costos reducidos $C_j - Z_j$.
    \item \textbf{Criterio de parada:} Se detiene el algoritmo cuando ningún costo reducido permite mejorar la función objetivo ($\le 0$ en maximización; $\ge 0$ en minimización). Los valores de las variables básicas se leen en la columna $b$ y las no básicas equivalen a cero.
\end{enumerate}

\section{Situaciones Especiales}

\begin{itemize}
    \item \textbf{Óptimos múltiples:} Ocurre si una variable no básica tiene costo reducido cero en la tabla óptima; puede entrar a la base sin modificar el valor óptimo de $Z$.
    \item \textbf{Degeneración:} Empate en la razón mínima que provoca que alguna variable básica tome valor cero, existiendo riesgo de ciclar sin avanzar.
    \item \textbf{No factibilidad:} La tabla final cumple la condición de óptimo pero mantiene una variable artificial básica con valor positivo, indicando incompatibilidad en las restricciones.
    \item \textbf{No acotación:} La columna de la variable entrante carece de coeficientes positivos; la variable puede incrementarse indefinidamente y la función objetivo no tiene cota finita.
\end{itemize}

\section{Comprobación Final y Aplicación Práctica}

\subsection*{Lista de verificación}
\begin{enumerate}
    \item Sustituir los valores obtenidos en las restricciones originales del problema.
    \item Confirmar el cumplimiento estricto de las condiciones de no negatividad y la coherencia de unidades.
    \item Recalcular de forma directa el valor de la función objetivo.
    \item Identificar e interpretar qué recursos se agotaron (holguras nulas) y cuáles quedaron con excedente.
\end{enumerate}

\subsection*{Soluciones continuas frente a enteras}
Si la decisión implica unidades indivisibles (como mobiliario o computadoras), la solución continua del método Simplex no debe redondearse al azar. Es necesario plantear un modelo de programación entera o analizar y justificar explícitamente cualquier redondeo para garantizar que ninguna restricción sea violada.

\end{document}